% === Appendix Table 1: Nearest vs All ===
\begin{table}[htbp]
  \centering
  \caption{\zihao{5} 最近球与所有球势能方案每回合球收集奖励均值对比}
  \label{tab:appendix_nearest_vs_all}
  \begin{tabular}{lcccc|c}
    \toprule
    方案 & 智能体0 & 智能体1 & 智能体2 & 智能体3 & 总平均 \\
    \midrule
    仅最近球 & 50.1 & 43.6 & 54.3 & 44.1 & 192.1 \\
    所有球 & 31.2 & 26.1 & 28.7 & 28.8 & 114.8 \\
    \midrule
    提升幅度 & \multicolumn{5}{c}{最近球方案总平均比所有球方案高出 67\%} \\
    \bottomrule
  \end{tabular}
\end{table}


% === Appendix Table 2: Potential Shaping v3 Comparison ===
\begin{table}[htbp]
  \centering
  \caption{\zihao{5} 不同奖励塑形方案每回合球收集奖励均值对比（v3，含所有方案）}
  \label{tab:appendix_potential_v3}
  \begin{tabular}{lcccc|c}
    \toprule
    方案 & 智能体0 & 智能体1 & 智能体2 & 智能体3 & 总平均 \\
    \midrule
    Linear & 100.5 & 89.0 & 91.4 & 98.4 & 379.4 \\
    Exponential & 81.2 & 76.0 & 77.8 & 75.0 & 310.1 \\
    Inverse & 100.6 & 93.4 & 91.7 & 94.6 & 380.2 \\
    Distance Reward &  \textbf{100.7} & 97.3 & 94.0 & 98.1 & 390.1 \\
    Sparse & 56.6 & 55.2 & 61.2 & 61.1 & 234.0 \\
    \midrule
    \multicolumn{6}{p{0.85\textwidth}}{%
      \small 注：五种方案在相同单智能体环境下训练 100 万步，4 个并行环境取均值。
      线性函数为 $\Phi(s)=-\frac{R_{\mathrm{ball}}}{r_{\mathrm{vision}}}d+R_{\mathrm{ball}}$，
      指数函数为 $\Phi(s)=R_{\mathrm{ball}}e^{-d/r_{\mathrm{vision}}}$，
      反比例函数为 $\Phi(s)=R_{\mathrm{ball}}\frac{r_{\mathrm{vision}}}{d+r_{\mathrm{vision}}}$，
      距离奖励为 $R=d_t-d_{t-1}$（靠近加分、远离减分），
      稀疏方案仅通过捡球得分学习。} \\
    \bottomrule
  \end{tabular}
\end{table}
